% explainer-source-hash: b65b60844a1e7b4a31378d5babc61291bbf8e105f40563de0308ddc7fe8ea1c6
% explainer-source-commit: b928cee68d8afd7dd077c24be1c4b2e0256110a9
% explainer-generated: 2026-07-16
% Regenerate with /explain-all (see WIRING.md). Do not edit this stamp by hand:
% it is how the toolchain knows whether this document still matches the code.
\documentclass[11pt,a4paper]{article}
\input{preamble}

\title{\textbf{PDF to Audiobook}\\[0.3em]\large A Brief}
\author{Portfolio explainer}
\date{}

\begin{document}
\maketitle
\thispagestyle{empty}

\section{The project in one page}

\begin{keyidea}{What it is}
A small desktop application, about 578 lines of Python in two files, that extracts the
text from PDF files and converts it into mp3 audiobooks. It handles a batch of files at
once, lets you preview the extracted text and restrict conversion to a page range,
offers a choice of language and regional accent, reports honest progress while it works,
and can play the result. It also runs entirely from the command line with no screen.
\end{keyidea}

The interesting thing about this project is not the feature list. It is that a hobby-scale
utility makes several decisions that larger codebases routinely get wrong, and that its
README is candid about the ones it did not solve.

\subsection{The three constraints that shaped it}

\begin{enumerate}
  \item \textbf{Speech synthesis happens on Google's servers.} The gTTS library is a thin
        client for Google Translate's speech endpoint. So the tool needs an internet
        connection, has no offline mode, is subject to rate limits, and has no volume
        control because the library offers none. No API key is needed;
        \file{.env.example} is deliberately empty and says so.
  \item \textbf{A PDF is not necessarily text.} A scanned PDF is a picture of a page.
        This project does no OCR, so such a document yields nothing.
  \item \textbf{A graphical interface freezes while your code runs.} Any slow work done
        directly in a button handler locks the window. Conversion is slow by nature. This
        single fact dictates the architecture of the GUI file.
\end{enumerate}

\section{Architecture}

The project is split in two, and the dependency points one way only.

\begin{figure}[H]
\centering
\resizebox{0.92\textwidth}{!}{
\begin{tikzpicture}[node distance=9mm]
  \node[box, minimum width=36mm] (gui) {\textbf{main.py} (360 lines)\\Tkinter GUI: batch queue,\\preview, options, progress};
  \node[box, minimum width=36mm, below=20mm of gui] (conv) {\textbf{converter.py} (218 lines)\\extraction, synthesis,\\playback, argparse CLI};
  \node[extbox, minimum width=26mm, left=20mm of conv] (cli) {Command line};
  \node[extbox, minimum width=20mm, below=18mm of conv] (gtts) {gTTS};
  \node[extbox, minimum width=20mm, left=12mm of gtts] (pypdf) {PyPDF2};
  \node[extbox, minimum width=24mm, right=12mm of gtts] (play) {ffplay / mpv /\\afplay / aplay};
  \node[store, minimum width=18mm, below=14mm of gtts] (goo) {Google TTS\\endpoint};
  \node[store, left=12mm of pypdf] (pdf) {PDF};
  \node[store, below=14mm of play] (mp3) {mp3};

  \draw[flow] (gui) -- node[lbl, right]{plain function calls;\\no GUI state crosses down} (conv);
  \draw[flow] (cli) -- (conv);
  \draw[flow] (conv) -- (pypdf);
  \draw[flow] (conv) -- (gtts);
  \draw[flow] (conv) -- (play);
  \draw[flow] (pypdf) -- (pdf);
  \draw[flow] (gtts) -- node[lbl, above]{HTTPS} (goo);
  \draw[flow] (play) -- (mp3);
\end{tikzpicture}}
\caption{\file{converter.py} never imports Tkinter and never touches a widget. That is
what makes both the command line entry point and headless verification possible.}
\end{figure}

\begin{keyidea}{Why the split is the load-bearing decision}
Because the logic module has no GUI dependency, it can be run and checked without a
screen. The README records that an earlier version was a single file mixing widgets and
conversion, and that separating them was a deliberate fix. Everything verified in
Section~\ref{sec:ver} of this document was verified by calling those functions directly,
which would have been impossible in the original design.
\end{keyidea}

\section{The two ideas worth understanding}

\subsection{Per-page synthesis is what makes progress real}

The obvious implementation sends the whole document to gTTS in one request. That request
is a black box: it takes a minute and reports nothing, so any progress bar drawn around
it is an animation on a timer, not information.

This project synthesizes \textbf{one page per request} instead. The code regains control
between pages, and that is where it calls an optional \code{progress\_callback(done,
total)}. The GUI turns that into ``Converting page 7/20''; the CLI prints \code{Page
7/20}; a script can pass nothing at all. The progress is truthful because the boundaries
are real. The cost, accepted deliberately, is more requests and more chances to hit a
rate limit.

\subsection{Joining the pages: raw mp3 concatenation}

Per-page synthesis leaves one mp3 fragment per page that must become one file. The code
keeps a single file handle open and writes each fragment's bytes into it, one after
another.

\begin{figure}[H]
\centering
\resizebox{0.9\textwidth}{!}{
\begin{tikzpicture}
  \node[box, minimum width=14mm, fill=accent!12] (a1) {frame};
  \node[box, minimum width=14mm, right=1mm of a1, fill=accent!12] (a2) {frame};
  \node[box, minimum width=14mm, right=6mm of a2, fill=okgreen!12] (b1) {frame};
  \node[box, minimum width=14mm, right=1mm of b1, fill=okgreen!12] (b2) {frame};
  \node[box, minimum width=14mm, right=6mm of b2, fill=warnred!10] (c1) {frame};
  \node[box, minimum width=14mm, right=1mm of c1, fill=warnred!10] (c2) {frame};
  \draw[decorate, decoration={brace, amplitude=4pt}, draw=accent]
    ([yshift=3mm]a1.north west) -- ([yshift=3mm]a2.north east) node[midway, above=3pt, lbl]{page 1};
  \draw[decorate, decoration={brace, amplitude=4pt}, draw=okgreen]
    ([yshift=3mm]b1.north west) -- ([yshift=3mm]b2.north east) node[midway, above=3pt, lbl]{page 2};
  \draw[decorate, decoration={brace, amplitude=4pt}, draw=warnred]
    ([yshift=3mm]c1.north west) -- ([yshift=3mm]c2.north east) node[midway, above=3pt, lbl]{page 3};
  \draw[decorate, decoration={brace, mirror, amplitude=5pt}, draw=inkblue]
    ([yshift=-3mm]a1.south west) -- ([yshift=-3mm]c2.south east)
    node[midway, below=5pt, lbl]{one file; a decoder reads frames left to right and never sees the joins};
\end{tikzpicture}}
\caption{An mp3 is a chain of self-describing frames, so gluing streams together works.}
\end{figure}

This would destroy most formats: concatenating two PNG files produces corruption, not a
taller image, because a PNG has one header for the whole picture. An mp3 is a stream of
independent frames, each with its own header, so a decoder simply carries on. The README
is honest that this leans on a property of the format rather than a documented guarantee,
and that a ``proper'' fix would add a dependency to perform the same byte copy with more
ceremony. It was checked here and it holds (Section~\ref{sec:ver}).

\section{How the interface stays responsive}

Tkinter has a hard rule: \textbf{widgets may only be touched from the thread that created
them}. So conversion cannot run on the main thread (the window would freeze) and cannot
update the widgets from the worker thread (undefined behaviour). The project resolves this
with the standard three-part pattern.

\begin{figure}[H]
\centering
\resizebox{\textwidth}{!}{
\begin{tikzpicture}
  \node[box, minimum width=28mm] (ui) at (0,0) {Main thread\\(Tk event loop)};
  \node[store, minimum width=22mm] (q) at (5.6,0) {\code{queue.Queue}};
  \node[box, minimum width=28mm] (w) at (11.2,0) {Worker thread};
  \node[extbox, minimum width=22mm] (c) at (16.4,0) {gTTS};
  \draw[dashed, draw=linegrey] (ui) -- (0,-6.2);
  \draw[dashed, draw=linegrey] (q) -- (5.6,-6.2);
  \draw[dashed, draw=linegrey] (w) -- (11.2,-6.2);
  \draw[dashed, draw=linegrey] (c) -- (16.4,-6.2);
  \draw[flow] (0,-1.0) -- node[lbl, above]{Convert All: disable button, start thread} (11.2,-1.0);
  \draw[flow] (11.2,-1.8) -- node[lbl, above]{one request per page} (16.4,-1.8);
  \draw[flow] (16.4,-2.6) -- node[lbl, above]{\code{progress\_callback(done, total)}} (11.2,-2.6);
  \draw[flow] (11.2,-3.4) -- node[lbl, above]{\code{put((...))}, never touches a widget} (5.6,-3.4);
  \draw[dflow] (0,-4.2) -- node[lbl, above]{\code{root.after(100, ...)} polls, drains with \code{get\_nowait()}} (5.6,-4.2);
  \draw[dflow] (5.6,-5.0) -- node[lbl, above]{messages} (0,-5.0);
  \node[lbl, anchor=west] at (0.2,-5.7) {only the main thread redraws the queue and the progress bar};
\end{tikzpicture}}
\caption{The queue gets its own lifeline because no arrow ever runs from the worker
straight to a widget. Every fact reaches the display through the queue.}
\end{figure}

Five message types flow through it: \code{file\_start}, \code{page\_progress},
\code{file\_done}, \code{file\_error} and \code{all\_done}. Each file is wrapped in its
own exception handler in the worker, so one bad PDF fails its own row and the batch
continues.

\section{Smaller decisions that show care}

\begin{itemize}
  \item \textbf{Atomic writes.} Output goes to \file{name\_tts.mp3.part} and is renamed
        into place with \code{os.replace} only on success, with a \code{try/finally}
        cleaning up. A conversion interrupted at page 7 of 20 leaves no half-file
        masquerading as a finished audiobook.
  \item \textbf{Graceful degradation.} \code{tkinterdnd2} is imported in a
        \code{try/except}. If it is missing, drag and drop is skipped, the window's
        subtitle says so plainly, and the app runs on through the file dialog.
  \item \textbf{No fake controls.} gTTS has no volume parameter, so there is no volume
        slider. The README says why. A dummy slider would have been easier and worse.
  \item \textbf{Honest failure from \code{play\_audio}.} It searches for \code{ffplay},
        \code{mpv}, \code{afplay} and \code{aplay} in order and returns a boolean, so the
        GUI can say ``no player found, here is the path'' instead of doing nothing. The
        README records that an early draft failed silently here.
  \item \textbf{\code{page.extract\_text() or ""}.} PyPDF2 returns \code{None} for a page
        it cannot read, which would crash the join. One idiom turns a crash into a blank
        page, and an all-blank range raises a clear error rather than writing an empty
        mp3.
  \item \textbf{Exact dependency pins} (\code{PyPDF2==3.0.1}, \code{gTTS==2.5.4},
        \code{tkinterdnd2==0.6.2}), which is the correct choice for an application.
\end{itemize}

\section{What was verified by running it}
\label{sec:ver}

These claims were not taken from the README on trust. They were re-checked by executing
the code against a generated 3-page PDF and a text-free PDF.

\begin{center}
\begin{tabular}{@{}p{5.4cm}p{9.2cm}@{}}
\toprule
\textbf{Claim} & \textbf{Result} \\
\midrule
Page-range extraction & Pages 2 to 3 returned exactly those pages, inclusive at both ends. \\
Empty-document handling & Raised the expected \code{ValueError}; no empty mp3 written. \\
Full conversion & Succeeded; progress fired 1/3, 2/3, 3/3. \\
\textbf{mp3 concatenation} & \code{file} reports a valid \code{MPEG ADTS, layer III} stream, and \code{ffmpeg} decoded the whole thing to WAV with no errors. The claim holds. \\
\textbf{Page range is real} & 3 pages produced \textbf{9.744 s} of audio; pages 2 to 3 produced \textbf{6.576 s}. An exact 3:2 ratio, 3.248 s per page, which is strong evidence the joins drop and duplicate nothing. \\
Player discovery & Found \code{ffplay} first, matching the intended order. \\
GUI freeze bug & Reproduced against a real Tk window. See below. \\
\bottomrule
\end{tabular}
\end{center}

\code{tkinterdnd2} was not installed on the verifying machine, so the drag-and-drop path
was not exercised; the documented fallback path is the one that ran.

\section{Honest findings}

\begin{honest}{The one serious bug: the interface can be permanently frozen by an ordinary click}
Pressing \textbf{Clear Queue} or \textbf{Remove Selected} while a conversion is running
bricks the window until restart. This was reproduced by running the real code, not
inferred.

\code{start\_conversion} disables only the Convert All button; the queue-editing buttons
stay live. Queue messages carry a \emph{list position} as the item's identity, and
clearing the queue invalidates every position in flight. The next poll raises
\code{IndexError}. The polling function's only \code{except} catches \code{queue.Empty},
so the error escapes, and execution never reaches the final line that re-arms the timer
via \code{root.after}. The self-rescheduling chain breaks and never restarts.

Observed: the \code{IndexError} fires, the queue never drains again, and Convert All
stays \code{disabled} even after \code{all\_done} is posted, because nothing is alive to
read it. The worker thread meanwhile finishes and writes its mp3s successfully, which the
user cannot learn from the frozen interface.

Either fix suffices: disable the queue-editing buttons during conversion, as Convert All
already is; or give messages a stable identity (the file path) instead of a list
position.
\end{honest}

\begin{honest}{The rest, worst first}
\begin{itemize}
  \item \textbf{An inverted or out-of-range page range blames the document.} Verified:
        requesting pages 3 to 1, or page 99 of a 3-page file, produces ``No extractable
        text found ... for the selected page range''. The document is full of text; the
        range is the problem. A user could conclude their PDF was a scan. The GUI escapes
        this only because its Save Page Range applies \code{min}/\code{max}; the CLI is
        fully exposed.
  \item \textbf{The CLI prints raw tracebacks at users.} Verified for a missing file and
        a bad range. A traceback is a developer's diagnostic, not a user message.
  \item \textbf{Silent overwrite.} The output filename derives only from the PDF's name,
        so converting pages 1 to 5 and then 6 to 10 of one file overwrites the first
        result without warning. For a tool whose headline feature is page ranges, that is
        a sharp edge.
  \item \textbf{Progress labels mislead on a range.} ``Converting page N/M'' counts
        positions within the selected range, so pages 5 to 7 display ``page 1/3''.
  \item \textbf{\code{QueueItem} swallows every exception,} reducing encrypted, corrupt
        and non-PDF files to an indistinguishable \code{?} with the real error discarded.
  \item \textbf{Language and accent are independent dropdowns,} so most of the 80
        offered combinations (Japanese with an Irish accent) are meaningless, and the
        accent list contains a duplicate: ``Default (.com)'' and ``US (.com)'' are the
        same value.
  \item \textbf{No automated tests.} There is no test directory. Every check above was
        written for this document and discarded. The clean module split means a pytest
        suite over \file{converter.py} would be easy; its absence is a choice, not an
        obstacle. The README lists this itself.
  \item \textbf{Offline it does nothing.} Correctly and prominently documented, but still
        the largest functional limit.
\end{itemize}
\end{honest}

\begin{keyidea}{The fair verdict}
The engineering instincts here are better than the project's size suggests: the one-way
module split, the atomic rename, the queue-based thread communication, the boolean-
returning player search, and the refusal to ship a fake volume slider are all decisions a
careful engineer makes deliberately. The README is unusually honest, including about
things this review independently confirmed.

The weaknesses cluster in one place: the seam between the GUI's data model and the worker
thread, where list positions are used as identities and only one of several dangerous
buttons is disabled during a conversion. That is where the only severe bug lives, and it
is a contained fix rather than a design flaw. Everything else on the list is polish.
\end{keyidea}

\section{Glossary}

\begin{description}[leftmargin=0pt, style=nextline]
  \item[Atomic rename] A filesystem operation that either completes fully or not at all,
        with no visible in-between state. Used here so a partial audiobook never appears
        under the final filename.
  \item[Callback] A function handed to another function so it can call back while
        working. Here the converter calls \code{progress\_callback} with a done count
        and a total after each page.
  \item[Daemon thread] A background thread that does not keep the program alive when the
        main window closes.
  \item[Event loop] The loop a GUI sits in, waiting for clicks and redrawing. While your
        code runs inside it, nothing redraws, which is why slow work must go elsewhere.
  \item[Graceful degradation] Losing an optional feature cleanly instead of crashing when
        something is unavailable.
  \item[mp3 frame] The unit an mp3 is built from. Each carries its own header, which is
        why concatenated mp3 streams still decode.
  \item[OCR] Optical Character Recognition: reading text out of a picture. This project
        has none, so scanned PDFs yield nothing.
  \item[\code{queue.Queue}] A container designed for safe hand-off between threads. The
        standard answer to inter-thread communication in Python.
  \item[Thread] A second line of execution inside one program, running concurrently.
  \item[TLD] Top-level domain, the end of a web address. gTTS uses it to pick a regional
        Google endpoint, which is how the ``accent'' setting works.
  \item[TTS] Text to Speech. Here it is a network service, not a local engine.
\end{description}

\end{document}
